\chapter{Introduction}

\section{Background}
\TemplateTodo{Explain the general context, ideal condition, real problem, research/practical gap, and why the study is necessary.}

Write the background coherently from the broader context toward the specific \DocTypeNameLower\ problem. Use verifiable scholarly sources such as journal articles, conference papers, books, standards, or official documents \cite{exampleReference}.

\TemplateTip{Avoid jumping into technical details before explaining the problem. Support important claims with citations.}

\section{Problem Statement}
\TemplateTodo{Convert the research problem into specific, measurable questions that can be answered through the methodology.}

\begin{enumerate}
  \item How can an approach be designed to address the main problem in the research object?
  \item How can the proposed approach be evaluated using appropriate metrics?
  \item What factors influence the research results based on the analysis?
\end{enumerate}

\section{Scope and Limitations}
\TemplateTodo{Define the scope so the study remains focused. The scope may cover data, method, location, period, variables, tools, or assumptions.}

\section{Research Objectives}
\TemplateTodo{Align objectives with the problem statements. Use verifiable verbs such as analyze, design, implement, evaluate, or compare.}

\section{Research Benefits}
\TemplateTodo{Explain theoretical and practical benefits. Avoid exaggerated claims; state realistic contributions for the scale of the \DocTypeNameLower.}
